\documentclass[11pt]{article}
\usepackage[margin=1in]{geometry}
\usepackage{amsmath,amssymb,amsthm}
\usepackage{booktabs,tabularx,array}
\usepackage{enumitem}
\usepackage[expansion=false,protrusion=false]{microtype}
\usepackage[hidelinks]{hyperref}
\usepackage{url}
\setlength{\emergencystretch}{2em}
\setlist[itemize]{nosep,leftmargin=*}
\setlist[enumerate]{itemsep=3pt,leftmargin=*}
\newtheorem{assumption}{Assumption}[section]
\newtheorem{definition}[assumption]{Definition}
\newtheorem{proposition}[assumption]{Proposition}
\newtheorem{corollary}[assumption]{Corollary}
\theoremstyle{remark}
\newtheorem{remark}[assumption]{Remark}
\newcommand{\R}{\mathbb R}
\newcommand{\X}{\mathcal X}
\newcommand{\T}{\mathcal T}
\newcommand{\dd}{\mathrm d}
\newcommand{\diag}{\operatorname{diag}}
\newcommand{\argmax}{\operatorname*{arg\,max}}
\newcommand{\argmin}{\operatorname*{arg\,min}}
\newcommand{\ind}{\mathrm{ind}}
\newcommand{\cen}{\mathrm C}
\newcommand{\net}{\mathrm{net}}

\title{Base to Based\\[3pt]
\large : Two-Graph Model of Fleet Coordination}
\author{Arya Chakrabarti, Arthur Chen}
\date{9/27/2026}

\begin{document}
\maketitle

\begin{abstract}
We model the battery fleet using two 
graphs.  They are coupled through a joint network with graph $G_1$ providing physical constraints and $G_2^i$ describing strategy. 

The motivation of this research was to prove existence and 
uniqueness of the profit gain in different strategies. We wanted to learn more how Base's profit stucture works, and now we're here.

We can model this as a game and find an an exact potential for the affine case that is distinct from both
competitive potential and fleet profit. From this, we can get a better understanding about the market making abilities of Base. 


\end{abstract}

\tableofcontents

\section{Scope}
\label{sec:setup}

A question is how much additional operating value a common owner
can obtain by jointly scheduling batteries, relative to an explicitly
specified independent controller. The two-graph construction separates
\emph{electrical feasibility} from \emph{ choice }.

Suppose we fix a finite horizon of dispatch intervals
$\T=\{0,\ldots,T-1\}$ with duration $\Delta>0$ hours. Let all state variables be
indexed by $t=0,\ldots,T$. 

\begin{table}[htbp]
\centering\small
\begin{tabularx}{\textwidth}{@{}lX@{}}
\toprule
Symbol & Meaning \\
\midrule
$i=1,\ldots,N$ & Battery index; $n(i)$ is its node in $G_1$ \\
$c_{it},g_{it}$ & Grid-side charging and discharging power, MW \\
$q_{it}=g_{it}-c_{it}$ & Battery net injection, MW; positive means discharge \\
$Q_t=\sum_iq_{it}$ & Fleet net injection, MW \\
$e_{it},E_i$ & Stored energy and usable energy capacity, MWh \\
$\eta_{c,i},\eta_{d,i}$ & Charging and discharging efficiencies in $(0,1]$ \\
$d_{nt}$ & Exogenous nodal net demand excluding controlled batteries, MW \\
$L_t$ & Exogenous system demand excluding the modeled fleet response, MW \\
$p_t(Q_t)=f_t(L_t-Q_t)$ & Counterfactual settlement price, \$/MWh \\
$a_t,\beta_t$ & Affine price intercept and impact slope: $p_t(Q_t)=a_t-\beta_tQ_t$ \\
$\kappa_i\ge0$ & Degradation charge per MWh of grid-side throughput \\
$x_i=(c_i,g_i,e_i)$ & Battery trajectory; $x=(x_1,\ldots,x_N)$ \\
$D_i(x_i)=\kappa_i\Delta\sum_t(c_{it}+g_{it})$ & Trajectory degradation cost, dollars \\
$\X_i,\X$ & Local and joint feasible trajectory sets \\
\bottomrule
\end{tabularx}
\caption{Variables and Meaning.}
\end{table}



\section{Graph 1: Electrical Constraints}
\label{sec:g1}

\subsection{Radial network and linearized flow}
Let $G_1=(\mathcal N_1,\mathcal E_1)$ be a rooted radial feeder with root $0$.
A fleet spanning several feeders can be represented by a forest, with an
independent root condition for each component. Each node is assigned the attributes of at least one battery, and multiple batteries may be assigned to one node.

For a directed edge $e=(m,n)$ away from the root, write
$\operatorname{down}(n)$ for the downstream node set. In a lossless
LinDistFlow approximation, real-power flow is
\begin{equation}
 P_{et}=\sum_{k\in\operatorname{down}(n)}d_{kt}
       -\sum_{i:\,n(i)\in\operatorname{down}(n)}q_{it}.
 \label{eq:flow}
\end{equation}
Let $R_{et}$ denote reactive-power flow in MVAr, avoiding confusion with fleet
injection $Q_t$. In the base model, reactive demands and $R_{et}$ are fixed,
and batteries inject or withdraw only real power. Let $r_e,x_e$ be
per-unit impedances on a common three-phase base $S_{\rm base}$ MVA, and
$v_{nt}=|V_{nt}|^2$ the per-unit squared voltage. Then
\begin{equation}
 v_{nt}=v_{0t}-\frac{2}{S_{\rm base}}
 \sum_{e\in\mathcal P(0\to n)}(r_eP_{et}+x_eR_{et}).
 \label{eq:voltage}
\end{equation}

\subsection{Thermal and voltage constraints}
For a specified apparent-power rating $\overline S_e$, impose
\begin{align}
 P_{et}^2+R_{et}^2&\le\overline S_e^2,
 &&e\in\mathcal E_1,\ t\in\T, \label{eq:thermal}\\
 \underline V_n^2\le v_{nt}&\le\overline V_n^2,
 &&n\in\mathcal N_1,\ t\in\T. \label{eq:voltagebounds}
\end{align}
When $R_{et}$ is fixed and $|R_{et}|\le\overline S_e$, define
$K_{et}=\sqrt{\overline S_e^2-R_{et}^2}$. Equation~\eqref{eq:thermal}
then becomes the linear constraint $-K_{et}\le P_{et}\le K_{et}$.
If $|R_{et}|>\overline S_e$, the prescribed operating point is infeasible.


\begin{definition}[Joint instantaneous feasibility]
Let $\mathcal A_t\subseteq\R^N$ be the set of battery injections satisfying
\eqref{eq:flow}--\eqref{eq:voltagebounds}, together with any specified
connection import/export limits. In the fixed-reactive model,
$\mathcal A_t$ is a closed convex polyhedron.
\end{definition}


\section{Graph 2: Battery Trajectories}
\label{sec:g2}

\subsection{Local dynamics and boundary conditions}
The local battery model is
\begin{align}
 e_{i,t+1}&=e_{it}+\Delta\left(\eta_{c,i}c_{it}
                        -\frac{g_{it}}{\eta_{d,i}}\right),
 &&t\in\T, \label{eq:soc}\\
 0\le c_{it}&\le\overline P_i^c,\qquad
 0\le g_{it}\le\overline P_i^d, &&t\in\T, \label{eq:powerboxes}\\
 \underline e_{it}&\le e_{it}\le E_i, &&t=0,\ldots,T, \label{eq:energybounds}\\
 e_{i0}&=e_i^{\rm init},\qquad e_{iT}=e_i^{\rm init}.
 \label{eq:boundary}
\end{align}
Here $\underline e_{it}$ is a declared household backup reserve.

The exact operating model additionally imposes
\begin{equation}
 c_{it}g_{it}=0,\qquad i=1,\ldots,N,\ t\in\T,
 \label{eq:exclusive}
\end{equation}

because a battery cannot be have net inflow and outflow of energy at the same time unless both values are 0. 

Let us have $\X_i^{\rm ex}$ denote this  nonconvex set. For convex
analysis, use the explicit relaxation $\X_i^{\rm rel}$ obtained by replacing
\eqref{eq:exclusive} with
\begin{equation}
 \frac{c_{it}}{\overline P_i^c}
 +\frac{g_{it}}{\overline P_i^d}\le1
 \label{eq:moderelax}
\end{equation}
assuming positive ratings. 

\begin{remark}[Why the relaxation certifies an optimum]
A relaxed optimum satisfying \eqref{eq:exclusive} is also a global optimum
of the exact problem with the same objective, since
$\X_i^{\rm ex}\subseteq\X_i^{\rm rel}$. 
\end{remark}

\subsection{Time-expanded graph representation}
For each battery, lets choose finite energy bins and a finite action menu.
$G_2^i$ has nodes $(t,e)$ and directed edges
$(t,e)\to(t+1,e')$ whenever a menu action $(c,g)$ satisfies
\eqref{eq:soc} and the local bounds. Fix the initial node and permitted
terminal nodes according to \eqref{eq:boundary}. A path specifies a battery
trajectory. At fixed prices, the profit attached to an edge is
\begin{equation}
 \Delta\big[p_t(g-c)-\kappa_i(c+g)\big].
 \label{eq:edgeprofit}
\end{equation}
Consequently, local price-taking optimization on this finite acyclic graph
is a longest-path problem, or a shortest-path problem after negating profits.


\subsection{Coupling the two graphs}
For $s\in\{\mathrm{ex},\mathrm{rel}\}$, define
\begin{equation}
 \X^s=\left\{x\in\prod_i\X_i^s:
                   (q_{1t},\ldots,q_{Nt})\in\mathcal A_t
                   \text{ for every }t\in\T\right\}.
 \label{eq:jointset}
\end{equation}
In the fixed-reactive model, write these shared constraints compactly as
$Ax\le b$. Each block $A_i$ maps battery $i$'s trajectory into the shared
network inequalities, so $Ax=\sum_iA_ix_i$.

$G_1$ excludes \emph{joint combinations} of transitions in the $G_2^i$; The joint strategy graph can be
viewed as a product graph with inadmissible combinations removed.

\begin{assumption}[]
\label{ass:feasible}
The horizon and all power/energy bounds are finite, and the chosen joint feasible set is nonempty. 
Network constraints and household reserves are identical across the regimes being compared.
\end{assumption}

\section{Price Model, Payoffs}
\label{sec:regimes}

\subsection{Aggregate price}
For each interval, simplify the price dependent on net injection being
\begin{equation}
 p_t(Q_t)=f_t(L_t-Q_t),
 \label{eq:price}
\end{equation}
where $f_t$ is continuous and nondecreasing on an interval containing $L_t$. 

The affine specialization is
\begin{equation}
 p_t(Q_t)=a_t-\beta_tQ_t,\qquad\beta_t\ge0,
 \label{eq:affine}
\end{equation}
with $\beta_t$ measured in $(\$/\mathrm{MWh})/\mathrm{MW}$.

At a fixed price vector $p$, battery $i$'s price-taking objective is
\begin{equation}
 U_i(x_i;p)=\sum_t\Delta p_tq_{it}-D_i(x_i).
 \label{eq:fixedpayoff}
\end{equation}
Its realized strategic payoff, when its action affects price, is instead
\begin{equation}
 \pi_i(x)=\sum_t\Delta q_{it}p_t(Q_t)-D_i(x_i).
 \label{eq:strategicpayoff}
\end{equation}

With a vector price $p(q)$, a competitive integral potential requires suitable
integrability conditions (in particular a symmetric Jacobian on a simply
connected differentiable domain).

\subsection{Threshold}
A threshold policy is a specified map from available information and local
state into charging, discharging, or idling, for example using thresholds
$\ell_{it}(e)<u_{it}(e)$. 

An independent horizon optimizer instead selects
\begin{equation}
 \widehat x_i\in\argmax_{x_i\in\X_i}U_i(x_i;\widehat p)
 \label{eq:independent}
\end{equation}
against a supplied  forecast $\widehat p$. Generally
$\widehat p\ne p(Q(\widehat x))$; therefore, independent optimization should not be
a competitive equilibrium either.

One explicit
offline benchmark with network protection is
\begin{equation}
 x^{\rm B}=\argmin_{x\in\X}
       \frac12(x-\widehat x)^\top W(x-\widehat x),\qquad W\succ0,
 \label{eq:projection}
\end{equation}
where the positive-definite, dimensionally scaled weights are fixed in advance.
On the convex set this projection exists and is unique.

\section{Competitive Equilibrium}
\label{sec:competitive}

\subsection{ Price-consistent competitive solution}

 Let $\X=\prod_i\X_i$.
 
A price-taking competitive solution is a profile $x^{\rm W}$ such that
\begin{equation}
 x_i^{\rm W}\in\argmax_{x_i\in\X_i}
       U_i\big(x_i;p(Q(x^{\rm W}))\big)
 \quad\text{for every }i.
 \label{eq:competitivefixedpoint}
\end{equation}
Each battery treats the price as fixed when deviating; the prices are then
required to be consistent with aggregate action. This is the competitive
price-taking condition, or its nonatomic/Wardrop. For a finite
fleet, it is not the Nash condition for the strategic payoff (which we will not prove here).
\eqref{eq:strategicpayoff}.

\subsection{Shared constraints}

Let $H_i$ be the linear map $x_i\mapsto q_i$. Define
\begin{equation}
 F_i^{\rm W}(x)=\nabla D_i(x_i)-\Delta H_i^\top p(Q(x)).
 \label{eq:competitivefield}
\end{equation}

\begin{definition}[Competitive variational equilibrium]
On the common convex feasible set $\X$, a competitive variational equilibrium
is $x^{\rm W}\in\X$ satisfying
\begin{equation}
 \sum_i\left\langle F_i^{\rm W}(x^{\rm W}),y_i-x_i^{\rm W}\right\rangle
 \ge0\qquad\text{for all }y\in\X.
 \label{eq:competitivevi}
\end{equation}
\end{definition}


Where Lagrange multipliers exist, a common network multiplier $\lambda\ge0$
supports individual optimization of
$U_i(x_i;p)-\lambda^\top A_ix_i$, with
$\lambda^\top(Ax-b)=0$.

\begin{proposition}[uniqueness]
\label{prop:potential}
Under Assumption~\ref{ass:feasible}, define
\begin{equation}
 \Phi(x)=\sum_t\Delta\int_0^{Q_t(x)}p_t(s)\,\dd s
                  -\sum_iD_i(x_i).
 \label{eq:potential}
\end{equation}
The function $\Phi$ is concave, its global maximizers exist and coincide
with the solutions of \eqref{eq:competitivevi}. On a product feasible set,
these are exactly the competitive solutions
\eqref{eq:competitivefixedpoint}. If each $p_t$ is strictly decreasing on
the feasible aggregate interval, all maximizers have the same aggregate
injection vector $Q^{\rm W}$. 
\end{proposition}

\begin{proof}
The derivative of $\int_0^{Q_t}p_t(s)\,\dd s$ is $p_t(Q_t)$, a
nonincreasing function, so the integral is concave. Since $Q$ is linear and
$D_i$ is convex, $\Phi$ is concave on $\X$; continuity and compactness give
existence. Its block gradient is
$\nabla_i\Phi=\Delta H_i^\top p(Q)-\nabla D_i=-F_i^{\rm W}$.
The first-order condition for a differentiable concave maximum on a convex
set is therefore exactly \eqref{eq:competitivevi}. On a product set this
condition is equivalent to the separate first-order conditions for each
concave fixed-price problem. If two maximizers had different aggregate
vectors, strict concavity of the integral in $Q$, combined with convexity of
the feasible set and costs, would make their midpoint strictly better, a
contradiction.
\end{proof}

If $f_t$ is differentiable, the second derivative of the integral with
respect to $Q_t$ is $-f_t'(L_t-Q_t)\le0$. Thus the convention in
\eqref{eq:potential} requires \emph{maximization}; a cost-minimization
convention must use $-\Phi$. The potential depends on full trajectories when
costs differ across batteries; aggregate $Q$ alone need not determine cost.
Strict aggregate uniqueness also does not imply uniqueness of a nonconvex
exact or discrete-path equilibrium.

\section{Finite-Player Nash Equilibrium}
\label{sec:nash}

With shared constraints, define
$\X_i(x_{-i})=\{y_i:(y_i,x_{-i})\in\X\}$ .

\begin{proposition}[Exact potential]
\label{prop:atomic}
Under \eqref{eq:affine}, the strategic game has exact payoff potential
\begin{equation}
 \Psi(x)=\sum_t\Delta\left[
       a_tQ_t-\frac{\beta_t}{2}Q_t^2
              -\frac{\beta_t}{2}\sum_iq_{it}^2\right]
              -\sum_iD_i(x_i).
 \label{eq:atomicpotential}
\end{equation}
Every global maximizer of $\Psi$ on a nonempty compact feasible set is a
GNE, and is an ordinary Nash equilibrium if the feasible set is a product.
On a convex product set, Nash equilibria are exactly the global maximizers.
On a shared convex set, global maximizers are exactly the
\emph{variational} Nash equilibria; other GNE need not maximize $\Psi$.
\end{proposition}

\begin{proof}
Fix $x_{-i}$. Write $Q_t=q_{it}+Q_{-i,t}$ and expand the terms of
\eqref{eq:atomicpotential} that depend on $x_i$. They are
\[
 \sum_t\Delta\big[a_tq_{it}-\beta_tq_{it}Q_{-i,t}
                         -\beta_tq_{it}^2\big]-D_i(x_i),
\]
which equal $\pi_i(x_i,x_{-i})$. All other terms are constant in $x_i$.
Hence any unilateral payoff difference equals the corresponding potential
difference, establishing the exact-potential identity. A global maximizer
therefore admits no profitable feasible unilateral deviation. Continuity and
compactness ensure that at least one such maximizer exists.
For convex sets, $\Psi$ is concave because $\beta_t\ge0$, and its block
gradient is the player's own-payoff gradient. On a product set, the separate
Nash first-order conditions combine into the global first-order condition
for $\Psi$; concavity makes it sufficient.
\end{proof}

The exact-potential identity also applies to a finite path-selection game
with affine impact. It guarantees existence of a pure equilibrium via a
global potential maximizer, but not that every pure equilibrium globally
maximizes the potential. For general nonlinear impact, the affine identity
does not extend without additional structure.


\section{Centralized Fleet Dispatch}
\label{sec:central}

\subsection{Existence}
Define fleet operating profit by
\begin{equation}
 \Pi(x)=\sum_i\pi_i(x)
       =\sum_t\Delta Q_t p_t(Q_t)-D(x),\qquad D=\sum_iD_i.
 \label{eq:fleetprofit}
\end{equation}
A centralized fleet optimum is
\begin{equation}
 x^{\cen}\in\argmax_{x\in\X}\Pi(x).
 \label{eq:central}
\end{equation}
Continuity and compactness ensure existence.

In the affine case,
\begin{equation}
 \Pi(x)=\sum_t\Delta(a_tQ_t-\beta_tQ_t^2)-D(x).
 \label{eq:centralqp}
\end{equation}
This is concave. 

For twice differentiable $f_t$, the revenue term
$h_t(Q)=Qf_t(L_t-Q)$ satisfies
\begin{equation}
 h_t''(Q)=-2f_t'(L_t-Q)+Qf_t''(L_t-Q).
 \label{eq:curvature}
\end{equation}
Thus $Qf_t''(L_t-Q)\le2f_t'(L_t-Q)$ on the feasible range is a
sufficient intervalwise condition for concavity. Convexity of $f_t$ alone
does not suffice. Piecewise curves and price caps require a separate
concavity assessment; a local optimum is not automatically global.

\subsection{Marginal Incentives}

The price-related marginal terms are
\begin{equation}
 \begin{array}{ll}
 \text{competitive price taker:} & p_t(Q_t),\\[2pt]
 \text{finite strategic battery }i: & p_t(Q_t)+q_{it}p_t'(Q_t),\\[2pt]
 \text{central fleet:} & p_t(Q_t)+Q_tp_t'(Q_t).
 \end{array}
 \label{eq:marginals}
\end{equation}
For affine impact, these become $p_t$, $p_t-\beta_tq_{it}$, and
$p_t-\beta_tQ_t$.

\subsection{Zero coordination value}
\begin{proposition}[Separability benchmark]
\label{prop:separable}
If prices are exogenous, 
\begin{equation}
 \max_{x\in\prod_i\X_i}\sum_iU_i(x_i;p)
       =\sum_i\max_{x_i\in\X_i}U_i(x_i;p).
 \label{eq:separable}
\end{equation}
This means centralized optimization cannot improve on independent optimal paths.
\end{proposition}
\begin{proof}
Every summand is bounded above by its individual maximum. Selecting an
individual maximizer for each battery attains the sum, since their product
is feasible. Compactness guarantees these maximizers exist.
\end{proof}



\section{Coordination Value: Valid Bounds}
\label{sec:premium}

\subsection{Common scoring and central dominance}
Let $x^{\rm B}\in\X$ be a specified feasible baseline: a protected threshold
policy, protected independent optimizer, competitive variational equilibrium,
or Nash equilibrium. 
\begin{equation}
 V_{\cen}=\Pi(x^{\cen}),\quad V_{\rm B}=\Pi(x^{\rm B}),\quad
 \Delta V_{\rm B}=V_{\cen}-V_{\rm B},\quad
 \rho_{\rm B}=\frac{V_{\cen}}{V_{\rm B}}\quad(V_{\rm B}>0).
 \label{eq:premiumdef}
\end{equation}

\begin{proposition}[]
\label{prop:dominance}
If $x^{\cen}$ is a global optimum and $x^{\rm B}$ is feasible for the same
problem, then $\Delta V_{\rm B}\ge0$. If also $V_{\rm B}>0$, then
$\rho_{\rm B}\ge1$. A feasible computed central schedule with value
$\underline V_{\cen}$ and a valid global upper bound $\overline V_{\cen}$
give
\[
 \underline V_{\cen}-V_{\rm B}\le\Delta V_{\rm B}
                    \le\overline V_{\cen}-V_{\rm B}.
\]
\end{proposition}


\subsection{Additive bound for affine impact}
\begin{proposition}[Market-impact premium bound]
\label{prop:impactbound}
Under affine impact, let $x^{\rm W}$ globally maximize the competitive
potential and $x^{\cen}$ globally maximize fleet profit on the same set.
Define
\begin{equation}
 K(x)=\frac{\Delta}{2}\sum_t\beta_tQ_t(x)^2.
 \label{eq:impactpenalty}
\end{equation}
Then
\begin{equation}
 0\le\Pi(x^{\cen})-\Pi(x^{\rm W})
 \le K(x^{\rm W})-K(x^{\cen})
 \le K(x^{\rm W}).
 \label{eq:impactbound}
\end{equation}
In particular $K(x^{\cen})\le K(x^{\rm W})$. If $\Pi(x^{\rm W})\ge v_0>0$
and $|Q_t|\le\overline Q_t$, then
\begin{equation}
 1\le\rho_{\rm W}
 \le1+\frac{\Delta\sum_t\beta_t\overline Q_t^2}{2v_0}.
 \label{eq:conditionalratio}
\end{equation}
\end{proposition}
\begin{proof}
For affine impact, $\Phi(x)=\sum_t\Delta(a_tQ_t-\beta_tQ_t^2/2)-D(x)$,
so $\Pi=\Phi-K$. Global potential optimality gives
$\Phi(x^{\cen})-\Phi(x^{\rm W})\le0$. Subtracting the corresponding
$K$ terms yields the middle bound; central optimality gives the lower
bound. Since $K\ge0$, the remaining inequalities follow.
\end{proof}


\subsection{No universal affine profit-ratio bound}
\begin{proposition}[Unbounded ratio in an affine two-period example]
\label{prop:counterexample}
There is no universal upper bound of $4/3$, or any other finite constant,
on $\rho_{\rm W}$.
\end{proposition}
\begin{proof}
Consider two one-hour intervals, lossless batteries, no degradation, no
binding feeder constraints, zero initial and terminal energy, and aggregate
cycling capacity $C<50$ MWh with matching power capacity. Each battery can
charge in period 1 and discharge in period 2. Let $z\in[0,C]$ be aggregate
energy shifted, so $Q_1=-z$ MW and $Q_2=z$ MW. Set
\[
 p_1=20+z,\qquad p_2=120-z.
\]
The competitive price spread is $100-2z>0$ for every $z\le C$, so each
price-taking battery uses its full capacity: $z^{\rm W}=C$. Fleet profit is
$\Pi(z)=100z-2z^2$. For $25\le C<50$, the centralized optimum is
$z^{\cen}=25$, with profit $1250$. Hence
\[
 \rho_{\rm W}(C)=\frac{1250}{100C-2C^2}\longrightarrow\infty
 \qquad\text{as }C\uparrow50.
\]
Every denominator in this limit is strictly positive, and all prices are
affine. This proves the claim.
\end{proof}

For $C=49$, the arithmetic is:
\begin{table}[htbp]
\centering\small
\begin{tabular}{@{}lrrrr@{}}
\toprule
Regime & Shifted MWh & Charge \$/MWh & Discharge \$/MWh & Profit (\$)\\
\midrule
Competitive & 49 & 69 & 71 & 98\\
Central fleet & 25 & 45 & 95 & 1,250\\
\bottomrule
\end{tabular}
\caption{Exact analytical example: $\rho_{\rm W}=1250/98\approx12.7551$..}
\end{table}

The same capacity-saturated profile is also a finite-player Nash equilibrium
for $N\ge49$ equally sized batteries in this $C=49$ example. At the proposed
profile, battery $i$'s own-profit derivative with respect to its shifted
energy is $100-2(49)-2(49/N)=2-98/N\ge0$. Its payoff is concave in that
energy, so its constrained optimum is its upper bound $49/N$.


\section{Fleet Profit}
\label{sec:welfare}

\subsection{System-cost}
Suppose, specifically for a welfare model, that $f_t(\ell)=C_t'(\ell)$, then
\begin{equation}
 C_t(L_t)-C_t(L_t-Q_t)=\int_0^{Q_t}f_t(L_t-s)\,\dd s.
 \label{eq:welfare}
\end{equation}
Consequently $\Phi(x)$ is avoided production cost minus degradation, and
its maximizer minimizes total system production cost plus degradation on
the same feasible set. Under these assumptions, competitive dispatch can be
system-efficient even when it earns the fleet less than \eqref{eq:central}.

Market settlement prices, scarcity adders, and offers need not equal true
marginal social costs. Without the additional assumption above, this welfare
interpretation is not justified. Throughout this paper, \eqref{eq:central}
is called \emph{fleet-profit optimal}, not a system social optimum. 

\subsection{Retail exposure changes the owner's objective}
For a simplified single-price retail portfolio, let $B_t$ be residual
unhedged household demand before controlled battery dispatch, with positive
$B_t$ denoting net purchases. Assume fixed retail receipts, fixed hedge
payments, and settlement of all residual energy at $p_t(Q_t)$.

The dispatch-dependent portfolio margin is
\begin{equation}
 \Pi^{\rm REP}(x)=\sum_t\Delta(Q_t-B_t)p_t(Q_t)-D(x)
                    +\text{constants}.
 \label{eq:rep}
\end{equation}
Relative to zero dispatch, its incremental value is
\begin{equation}
 V^{\rm REP}(x)=\sum_t\Delta\left[(Q_t-B_t)p_t(Q_t)+B_tp_t(0)\right]-D(x).
 \label{eq:repvalue}
\end{equation}

For affine impact, the marginal price component is
\begin{equation}
 p_t(Q_t)-\beta_t(Q_t-B_t).
 \label{eq:repmarginal}
\end{equation}
Lower prices benefit the owner's remaining net purchases. Therefore
$B_t$ can change optimal dispatch. 

Base's published customer documentation states that it owns and schedules
its batteries \cite{base}. Accordingly, independent controllers are a
counterfactual operating architecture.

\section{Coincident-Peak Timing (4CP)}
\label{sec:cp}

How we understand it, four-coincident-peak (4CP) transmission-cost regime is the institutional
reason a single 15-minute interval can carry outsized economic weight in
ERCOT. 

We use it in two ways: as \emph{context} that motivates treating the
coincident-peak (CP) interval as a distinguished target on $G_2$, and as a
stochastic extension of the strategy graph. 

\subsection{The 4CP mechanism and its scheduled replacement}
Under 4CP, ERCOT allocates transmission cost of service to distribution
providers according to each provider's load during the single highest
15-minute system-wide demand interval in each of June, July, August, and
September; the following calendar year's transmission charge scales with that
coincident-peak demand \cite{ercot4cp,rtoinsider}. 

This policy is scheduled to change:
under Texas Senate Bill 6, the Public Utility Commission of Texas must overhaul
the methodology by December 31, 2026, with a twelve-coincident-peak (12CP)
approach measured in 30-minute intervals under active consideration in Project
No.~58484 \cite{sb6reform,rtoinsider}. 

We therefore parameterize the setting by
a target-interval count $m$ and an interval length $\Delta_{\rm CP}$, rather
than fixing $m=4$ and $\Delta_{\rm CP}=15$ minutes.

\subsection{Who is exposed}
4CP is a billing determinant only for interval/demand-metered accounts, in
practice larger commercial and industrial loads. Residential and
small-commercial transmission cost is recovered through volumetric
(\$/kWh) or fixed allocations that are not tied to the coincident-peak
intervals; there is accordingly no direct per-meter transmission saving from
residential curtailment during a CP.

\subsection{The coincident-peak interval as an uncertain target}
Within a CP window, let the realized peak interval $\tau$ be a random element
of a candidate set $\mathcal C\subseteq\T$ with distribution $\pi$ induced by
weather, demand, renewable output, and by the aggregate
curtailment of all CP-responsive load. Aggregate avoidance flattens and can
displace the peak; this signal degradation is documented and is part of the
regulatory case for reform \cite{nrg,rtoinsider}. A fleet that is a negligible
share of system load does not materially move $\pi$ and best-responds to
it. 

\subsection{Deployment as optimal stopping under a shared budget}
Over the window, several candidate high-load intervals are revealed
sequentially with forecasts. The fleet holds a limited, costly discharge
capability: cycle wear, energy conversion loss, and household backup reserve
$\underline e_{it}$ that is most valuable precisely on the hottest,
highest-outage-risk days. Choosing when to commit the fleet against the running
maximum is an optimal-stopping (prophet-type) problem; committing on a false
alarm consumes budget and backup that may be needed for the true peak. This is
the intertemporal counterpart, under timing uncertainty, of the deterministic
dispatch in Sections~\ref{sec:central}--\ref{sec:premium}.

\subsection{When does coordinating the deployment add value?}
The separability logic of Proposition~\ref{prop:separable} transfers directly.

\begin{proposition}[Zero deployment-coordination premium under separability]
\label{prop:cpseparable}
If the timing distribution $\pi$ is exogenous, each battery has its own
budget, the feasible set is a product, and the objective is additive over
batteries and risk-neutral, then the fleet-optimal expected deployment equals
the sum of independently optimal per-battery deployments, and centralized
deployment adds no expected value.
\end{proposition}
\begin{proof}
Under the stated hypotheses the expected objective is a sum of per-battery
functionals over per-battery feasible sets, so the maximization separates
exactly as in Proposition~\ref{prop:separable}; the expectation over $\pi$,
common to all batteries, does not couple them.
\end{proof}


\subsection{Placement in the two-graph model}
The CP setting adds a stochastic weight to $G_2$: a trajectory that is armed
and available at the realized $\tau$ carries an option value proportional to
the monetizable peak spread in \eqref{eq:repvalue}, while $G_1$ supplies the
shared constraint of channel~1 that can make coordinated deployment strictly
better than independent deployment. Operationally, per-customer (micro) prediction
supplies each battery's availability and cost inputs, and the fleet (macro)
layer solves the stochastic, feeder-constrained, budget-limited deployment
against $\pi$.

\section{Large-Fleet Scaling Without an Unsupported Limit Claim}
\label{sec:limit}

Consider a sequence of fleets indexed by $N$, and write
$q_{it}^{(N)}=a_{it}^{(N)}/N$. Then
\begin{equation}
 Q_t^{(N)}=\sum_iq_{it}^{(N)}
           =\frac1N\sum_i a_{it}^{(N)}.
 \label{eq:scaling}
\end{equation}
If the empirical distribution of the rescaled actions converges weakly to
$\mu_t$ and these actions are uniformly bounded, then
$Q_t^{(N)}\to\int a\,\mu_t(\dd a)$. By contrast,
$N^{-1}\sum_iq_{it}^{(N)}=Q_t^{(N)}/N\to0$.


For affine impact the finite-player and competitive potentials satisfy
\begin{equation}
 \Psi(x)=\Phi(x)-\frac{\Delta}{2}\sum_t\beta_t\sum_iq_{it}^2.
 \label{eq:limitpotential}
\end{equation}
If $|q_{it}|\le\varepsilon_N\to0$ and
$\sum_i|q_{it}|\le M_t$ uniformly, then
\begin{equation}
 0\le\Phi(x)-\Psi(x)
 \le\frac{\Delta\varepsilon_N}{2}\sum_t\beta_tM_t\longrightarrow0.
 \label{eq:uniformbound}
\end{equation}
The bound follows from $q_{it}^2\le\varepsilon_N|q_{it}|$. 

It shows why
individual price-impact corrections vanish while the fleet correction
$\beta_tQ_t$ can remain finite. 


\section{Numerical Interpretation and Verification Protocol}
\label{sec:numerics}

\subsection{Previously reported results retained for traceability}
The original numerical study reported ERCOT real-time 15-minute
\texttt{HB\_HUBAVG} prices for an unspecified shoulder-season day, with
minimum/mean/maximum of approximately \$18/40/101 per MWh, an inverse-supply
model with a scarcity component, and 5 kW/15 kWh units with round-trip
efficiency 0.81. Table~\ref{tab:reported} preserves its reported dollar values.


\begin{table}[htbp]
\centering\small
\begin{tabular}{@{}lrrrr@{}}
\toprule
Scenario & Units & Baseline (\$) & Central (\$) & $\Delta V$ (\$)\\
\midrule
Calm & 50,000 & 11,682 & 11,683 & 1\\
Calm & 150,000 & 34,876 & 34,884 & 8\\
Scarcity & 50,000 & 144,779 & 144,781 & 2\\
Scarcity & 150,000 & 323,061 & $\sim325,000$ & $\sim1,939$\\
\bottomrule
\end{tabular}

\label{tab:reported}
\end{table}


\subsection{ERCOT anchoring}

One transparent counterfactual specification is
\begin{equation}
 p_t(Q_t)=p_t^{\rm obs}-\beta_t(Q_t-Q_t^{\rm ref}).
 \label{eq:historicalanchor}
\end{equation}
The reference dispatch $Q_t^{\rm ref}$ must match the price anchor.



\section{What the Two-Graph Model Establishes}
The physical graph supplies shared feasibility constraints; the strategy
graphs supply local intertemporal choices. Their combination gives a
well-defined dispatch problem without identifying synchronization with
equilibrium or assuming that centralization must create substantial value.
Competitive, finite strategic, and fleet-optimal dispatch have different
objectives. Their relationship can be proved under stated convexity and
price assumptions, and evaluated against the same physical constraints.

The most direct empirical case for centralized control is a measured
improvement over a declared feasible independent baseline when feeder limits
or other shared constraints matter. Market-impact value is a distinct,
assumption-sensitive component. Coincident-peak timing, driven by the
4CP regime and its scheduled 12CP successor, sharpens \emph{when} that value
concentrates but, for a residential fleet, is realized through the energy and
portfolio objectives rather than a transmission credit. The retained
simulations can support claims about their own scenarios after verification,
while the propositions above specify which conclusions hold beyond those
scenarios.

\appendix


\begin{thebibliography}{99}
\bibitem{rosenthal}
R. W. Rosenthal, ``A class of games possessing pure-strategy Nash
equilibria,'' \emph{International Journal of Game Theory}, 2, 65--67, 1973.
\url{https://doi.org/10.1007/BF01737559}.

\bibitem{roughgarden}
T. Roughgarden and E. Tardos, ``How bad is selfish routing?''
\emph{Journal of the ACM}, 49(2), 236--259, 2002.

\bibitem{base}
Base Power, ``Can I control when the Base battery charges and discharges?''
Customer help documentation, edited March 19, 2026.
\url{https://help.basepowercompany.com/en/articles/10282753}.

\bibitem{ercot4cp}
Public Utility Commission of Texas, Project No.~58484, ``Evaluation of
Transmission Cost Recovery Allocation,'' staff questions and stakeholder
responses on the current four-coincident-peak (4CP) methodology and the
referenced ERCOT postage-stamp transmission rate, 2025.
\url{https://interchange.puc.texas.gov/search/documents/?controlNumber=58484}.

\bibitem{nrg}
NRG Energy, ``Comments on Staff Questions, Project No.~58484,'' Sept.~9, 2025
(residential/small-commercial volumetric transmission recovery; divergence of
4CP from scarcity pricing).
\url{https://www.nrg.com/assets/documents/energy-policy/nrg-comments-on-transmission-cost-recovery-puct-staff-rfc-1-090925.pdf}.

\bibitem{huntsman4cp}
K. Van Horn et al., ``The response of large industrial energy consumers to
four coincident peak (4CP) transmission charges in the Texas (ERCOT) market,''
\emph{Electricity Journal}, 2013 (no direct 4CP benefit to an individual
residential or small-commercial consumer under current rate design).

\bibitem{sb6reform}
Texas Senate Bill 6 (2025) and PUCT rulemaking, Project No.~58484: mandate to
overhaul transmission cost allocation, with a 12-coincident-peak, 30-minute
proposal, target implementation by December 31, 2026.

\bibitem{rtoinsider}
RTO Insider, ``ERCOT's 4CP Mechanism Living on Borrowed Time,'' 2025
(PUCT development of a 12CP methodology measured in 30-minute intervals).
\url{https://www.rtoinsider.com/132246-ercot-4cp-mechanism-living-on-borrowed-time/}.

\bibitem{peak2026}
ERCOT summer resource assessment, 2026 (projected summer peak near 95 GW),
as reported in contemporaneous 4CP-season market commentary.

\bibitem{ercot-lmp}
ERCOT, ``LMPs by Resource Nodes, Load Zones and Trading Hubs,''
data product NP6-788-CD.
\url{https://www.ercot.com/mp/data-products/data-product-details?id=NP6-788-CD}.

\bibitem{ercot-spp}
ERCOT, ``Settlement Point Prices at Resource Nodes, Hubs and Load Zones,''
data product NP6-905-CD.
\url{https://www.ercot.com/mp/data-products/data-product-details?id=NP6-905-CD}.

\bibitem{ercot-history}
ERCOT, ``Historical RTM Load Zone and Hub Prices,''
data product NP6-785-ER.
\url{https://www.ercot.com/mp/data-products/data-product-details?id=NP6-785-ER}.

\bibitem{ercot-rtcb}
ERCOT, ``ERCOT Goes Live with Real-Time Co-optimization Plus Batteries
(RTC+B),'' December 5, 2025.
\url{https://www.ercot.com/news/release/12052025-ercot-goes-live}.

\bibitem{ercot-ordc}
ERCOT, \emph{2024 Biennial ERCOT Report on the ORDC}, October 31, 2024.
\url{https://www.ercot.com/files/docs/2024/10/31/2024-biennial-ercot-report-on-the-ordc-20241031.pdf}.
\end{thebibliography}

\end{document}
